\section{5. Discusión}
Los resultados muestran que el crecimiento agregado puede ocultar trayectorias distintas entre empresas. En industria, el indicador mejora porque cambia la composición del empleo, aunque el cambio promedio dentro de las empresas que permanecen sea negativo. En servicios recientes, ambos márgenes contribuyen al aumento. El comercio presenta señales positivas en 2020--2024, pero su ventana y sus identificadores limitan la comparación temporal.

La literatura ofrece ejemplos de los dos patrones. México y el estudio del Banco Mundial para Colombia destacan la debilidad del cambio dentro de las empresas; Chile y el Reino Unido muestran aportes importantes de reasignación. Portugal ilustra que las entrantes pueden reducir el promedio por diferencias de tamaño. Estas comparaciones ayudan a interpretar nuestros resultados, pero sus medidas y muestras no son idénticas.

El ejercicio permite describir de dónde proviene el cambio medido en las encuestas. Para establecer por qué ocurrió y qué implicaciones tiene, el siguiente análisis debe armonizar la muestra industrial con el estudio del Banco Mundial y comprobar la sensibilidad a deflactores, cobertura y definiciones de entrada y salida.
